\documentclass[a4paper,14pt]{extarticle}

\usepackage[utf8]{inputenc}
\usepackage[russian]{babel}
\usepackage[a4paper, margin=2.5cm]{geometry}

\usepackage{amssymb}
\usepackage{amsmath}

\usepackage{setspace}

\onehalfspacing

\begin{document}
    \begin{center}
        \textbf{Игра в монету}
    \end{center}

    Вероятность выпадания орла: $p = 0.5 = P(A)$

    Вероятность выпадания решки: $q = 1 - p = 0.5$

    Пусть $P_i, i = \overline{1,4}$ -- вероятности победы $i$ игрока на первом круге.\\
    Если выиграет первый игрок, то это значит, что ему выпал орёл: $P_1 = P(A) = p = 0.5$\\
    Если выиграет второй игрок, то это значит, что первому игроку выпала решка, а второму игроку -- орёл: $P_2 = P(A)P(\overline{A}) = pq = p^2 = 0.25$\\
    Если выиграет третий игрок, то это значит, что первому и второму игрокам выпала решка, а третьему -- орёл: $P_3 = P(A)P(\overline{A})^2 = pq^2 = p^3 = 0.125$\\
    Если выиграет четвёртый игрок, то это значит, что первому, второму и третьему игрокам выпала решка, а четвертому -- орёл: $P_4 = P(A)P(\overline{A})^3 = pq^3 = p^4 = 0.0625$\\

    Однако может случится так, что четвёртому игроку тоже не выпадет орёл, в таком случае игра продолжается дальше и монету бросает первый игрок. Первый игрок выиграет во втором кругу, если у всех игроков выпала решка, а первому игроку с второй попытки -- орёл.
    Обозначим вероятность этого события как $P_5 = P(A)P(\overline{A})^4 = pq^4q = p^5 = 0.03125$.

    Таким образом, продолжая так до тех пор, пока кому-то не выпадет орёл, получили формулу вероятности выпадения орла на $k$-м броске монеты:\\
    $P_k = P(A)P(\overline{A})^{k - 1} = pq^{k - 1} = p^k$

    Пусть события:
    \begin{itemize}
        \item $A_1$ -- первый игрок выиграл;
        \item $A_2$ -- второй игрок выиграл;
        \item $A_3$ -- третий игрок выиграл;
        \item $A_4$ -- четвёртый игрок выиграл.
    \end{itemize}

    Тогда получаем формулы вероятностей победы каждого игрока:\\
    $\displaystyle
    P(A_1) = \sum_{i = 0}^{\infty} P_{1 + 4i} = \sum_{i = 0}^{\infty} p^{1 + 4i} = \frac{p}{1 - p^4} = \frac{0.5}{1 - 0.5^4} = \frac{8}{15}\\
    P(A_2) = \sum_{i = 0}^{\infty} P_{2 + 4i} = \sum_{i = 0}^{\infty} p^{2 + 4i} = \frac{p^2}{1 - p^4} = \frac{0.25}{1 - 0.5^4} = \frac{4}{15}\\
    P(A_3) = \sum_{i = 0}^{\infty} P_{3 + 4i} = \sum_{i = 0}^{\infty} p^{3 + 4i} = \frac{p^3}{1 - p^4} = \frac{0.125}{1 - 0.5^4} = \frac{2}{15}\\
    P(A_4) = \sum_{i = 0}^{\infty} P_{4 + 4i} = \sum_{i = 0}^{\infty} p^{4 + 4i} = \frac{p^4}{1 - p^4} = \frac{0.0625}{1 - 0.5^4} = \frac{1}{15}
    $

\end{document}